\documentclass[12pt]{article}
\usepackage[utf8]{inputenc} % default from sharelatex
\usepackage[a4paper, left=30mm, right=30mm, top=30mm, bottom=30mm]{geometry}
% dont break my code
\usepackage{indentfirst} % to indent fist paragraph
\usepackage[brazilian]{babel} % BR
\usepackage{listings} % to show code
\usepackage{xcolor} % to colors
\usepackage{hyperref} % to links
\usepackage{graphicx} % to have pictures
\usepackage{float} % my figure, my place

\graphicspath{{svg/}}

\lstset{
  language=Prolog,
  basicstyle=\ttfamily\small,
  numberstyle=\footnotesize,
  numbers=left,
  backgroundcolor=\color{gray!10},
  frame=single,
  tabsize=2,
  rulecolor=\color{black!30},
  title=\lstname,
  escapeinside={\%*}{*)},
  breaklines=true,
  breakatwhitespace=true,
  framextopmargin=2pt,
  framexbottommargin=2pt,
  extendedchars=false,
  inputencoding=utf8,
  commentstyle=\color{blue},
  keywordstyle=\color{purple},
  showstringspaces=false,
  stringstyle=\color{red}
}

\hypersetup{
  colorlinks=true,
  urlcolor=blue
}

\title{t2C}

\author{Gustavo Borges França (15100741)
        \and
        Lucas João Martins (15100752)}
\date{}

\begin{document}

\maketitle

\section{Código dos testes}
\lstinputlisting{testes.pl}

\section{Imagens}
Essa foi a imagem utilizada antes da execução de qualquer um dos testes:
\begin{figure}[H]
  \includegraphics[width=\linewidth]{antes_testes.png}
  \caption{Figura base}
\end{figure}

Após a execução do \emph{testefigclone} essa é a imagem resultante:
\begin{figure}[H]
  \includegraphics[width=\linewidth]{01_depois_testefigclone.png}
  \caption{Figura após testefigclone}
\end{figure}

A partir da figura base e após a execução do \emph{testefigpf} a imagem resultante é:
\begin{figure}[H]
  \includegraphics[width=\linewidth]{02_depois_testefigpf.png}
  \caption{Figura após testefigpf}
\end{figure}

Novamente a partir da figura base e após a execução do \emph{testefigpt} a imgem resultante é:
\begin{figure}[H]
  \includegraphics[width=\linewidth]{03_depois_testefigpt.png}
  \caption{Figura após testefigpt}
\end{figure}

Mais uma vez a partir da figura base e após a execução do \emph{testefiggd} a imagem resultante é:
% \begin{figure}[h]
%   \includegraphics[width=\linewidth]{04_depois_testefiggd.png}
%   \caption{Figura após testefiggd}
% \end{figure}

Por fim, a partir também da figura base e após a execução do \emph{testefigge} a imagem resultante é:
% \begin{figure}[h]
%   \includegraphics[width=\linewidth]{05_depois_testefigge.png}
%   \caption{Figura após testefigge}
% \end{figure}

\section{Código-fonte}
\href{https://github.com/lucasjoao/paradigmas/tree/master/Prolog/T2c}
{GitHub}
\end{document}
